\documentclass[11pt]{article}
\usepackage[margin=1in]{geometry}
\usepackage{amsmath}
\usepackage{setspace}

\begin{document}
\onehalfspacing

\section*{Shadow Price Challenge}

Consider the linear program
\[
\max_{x_1,x_2}\quad z=\pi_1x_1+\pi_2x_2
\]
subject to
\[
\begin{aligned}
a_{11}x_1+a_{12}x_2 &\le b_1\\
a_{21}x_1+a_{22}x_2 &\le b_2\\
a_{31}x_1+a_{32}x_2 &\le b_3\\
x_1,x_2&\ge 0.
\end{aligned}
\]Recall that the shadow price of resource \(j\) is the marginal change in the optimized objective value from a small increase in \(b_j\):
\[
\lambda_j=\frac{\partial z^*}{\partial b_j}.
\]

Sam proposes the following shortcut:
\[
\boxed{\lambda_j=\frac{\pi_i}{a_{ji}}}
\]

That is, “the shadow price of a constraint is just the objective-function coefficient of a variable divided by the amount of that resource required to produce one unit of the variable.”



\textbf{Your task: Determine whether this statement is always true.}


If it is not always true, identify a set of conditions under which it is true. Your answer should include a mathematical argument, not just a numerical example.

As you work, consider the following questions:
\begin{itemize}
\item What must be true about constraint \(j\) at the optimum for it to have a nonzero shadow price?
\item If \(b_j\) increases slightly, which decision variable or variables change?
\item What would have to be true for
\[
  \frac{\partial x_i}{\partial b_j}=\frac{1}{a_{ji}}?
  \]
\item What could cause the shortcut to fail even if constraint \(j\) is binding?
\item Does your argument apply to arbitrarily large changes in \(b_j\), or only to sufficiently small changes?
\end{itemize}
Challenge: Draw a graph of a two-variable LP for which the shortcut works, and then modify the problem so that the same shortcut fails.

\end{document}